% GENERATED FILE. Edit canonical lessons, then run npm run generate:books.
% canonical-source-hash: fnv1a64:95c1f533c713268f
% canonical-lessons: GU-C206-pahocvu, GU-C206-thakvu, GU-C206-pravasi, GU-C206-mukam, GU-C206-gati, GU-R206-review-201-206

\chapter{Arriving at last}
\label{ch:gu-arriving-at-last}
\hlchaptermodality{206}

\begin{chapteropening}
\textbf{By the end of this chapter:} \emph{I can say when we arrived, how tired I was and how fast we went.}

Complete the last lesson of chapter 206: I can say when we arrived, how tired I was and how fast we went.
\end{chapteropening}

\section[\texorpdfstring{\gu{પહોંચવું} (pahõcvũ)}{પહોંચવું (pahõcvũ)}]{\gu{પહોંચવું} (pahõcvũ) — to arrive, to reach}

\label{lesson:GU-C206-pahocvu}



\begin{quote}
Before the new one: say the Gujarati for to halt, then the Gujarati for to get off.
\end{quote}

\subsection*{You'll want to know: \gu{પહોંચવું}}
\textbf{\gu{પહોંચવું}} — \emph{pahõcvũ} — \textquotedblleft{}to arrive, to reach\textquotedblright{}.

\begin{quote}
\textbf{\gu{અમે} \gu{રાત્રે} \gu{શહેરમાં} \gu{પહોંચ્યા}.} — \emph{ame rātre śahermā̃ pahõcyā} — \textquotedblleft{}We reached the city at night.\textquotedblright{}
\end{quote}

\begin{cousinweb}
\textbf{\gu{પહોંચવું}} (\emph{pahõcvũ}) is a tadbhava, a word inherited through Prakrit and reshaped by sound change. CDIAL lists it under Sanskrit \emph{prábhūta}, \textquotedblleft{}much, great\textquotedblright{}, by way of Prakrit \emph{pahuccaï}, \textquotedblleft{}is powerful, reaches\textquotedblright{} (CDIAL 8716).
\end{cousinweb}

\begin{grammarlens}[title={What you've built}]
One more word for everyday life at A2.
\end{grammarlens}

\subsection*{Your turn}
\begin{itemize}
\raggedright
  \item \emph{Say it:} \emph{pahõcvũ}
  \item \emph{Say it:} \emph{pahõcvũ}, once more
  \item \emph{Say it:} \emph{ame rātre śahermā̃ pahõcyā} — We reached the city at night.
  \item \emph{Recall:} say the Gujarati for to halt, then the Gujarati for to get off, then say \emph{pahõcvũ} again
\end{itemize}

\begin{tcolorbox}[breakable,colback=teal!4,colframe=teal!35!black,title={Before you move on}]
What does \gu{પહોંચવું} mean? (\textquotedblleft{}To arrive, to reach\textquotedblright{}.) Say it once more.
\end{tcolorbox}
\section[\texorpdfstring{\gu{થાકવું} (thākvũ)}{થાકવું (thākvũ)}]{\gu{થાકવું} (thākvũ) — to get tired}

\label{lesson:GU-C206-thakvu}



\begin{quote}
Before the new one: say the Gujarati for to get off, then the Gujarati for to arrive.
\end{quote}

\subsection*{You'll want to know: \gu{થાકવું}}
\textbf{\gu{થાકવું}} — \emph{thākvũ} — \textquotedblleft{}to get tired\textquotedblright{}.

\begin{quote}
\textbf{\gu{હું} \gu{લાંબી} \gu{મુસાફરીથી} \gu{થાક્યો}.} — \emph{hũ lā̃bī musāpharīthī thākyo} — \textquotedblleft{}I got tired from the long journey.\textquotedblright{}
\end{quote}

\begin{cousinweb}
\textbf{\gu{થાકવું}} (\emph{thākvũ}) is a tadbhava. It comes from \emph{*sthakk}, \textquotedblleft{}stop, halt\textquotedblright{}, a reconstructed form; CDIAL glosses the Gujarati verb \textquotedblleft{}to become weary\textquotedblright{} (CDIAL 13737).
\end{cousinweb}

\begin{grammarlens}[title={What you've built}]
One more word for everyday life at A2.
\end{grammarlens}

\subsection*{Your turn}
\begin{itemize}
\raggedright
  \item \emph{Say it:} \emph{thākvũ}
  \item \emph{Say it:} \emph{thākvũ}, once more
  \item \emph{Say it:} \emph{hũ lā̃bī musāpharīthī thākyo} — I got tired from the long journey.
  \item \emph{Recall:} say the Gujarati for to get off, then the Gujarati for to arrive, then say \emph{thākvũ} again
\end{itemize}

\begin{tcolorbox}[breakable,colback=teal!4,colframe=teal!35!black,title={Before you move on}]
What does \gu{થાકવું} mean? (\textquotedblleft{}To get tired\textquotedblright{}.) Say it once more.
\end{tcolorbox}
\section[\texorpdfstring{\gu{પ્રવાસી} (pravāsī)}{પ્રવાસી (pravāsī)}]{\gu{પ્રવાસી} (pravāsī) — tourist, traveller}

\label{lesson:GU-C206-pravasi}



\begin{quote}
Before the new one: say the Gujarati for to arrive, then the Gujarati for to get tired.
\end{quote}

\subsection*{You'll want to know: \gu{પ્રવાસી}}
\textbf{\gu{પ્રવાસી}} — \emph{pravāsī} — \textquotedblleft{}tourist, traveller\textquotedblright{}.

\begin{quote}
\textbf{\gu{ગામમાં} \gu{ઘણા} \gu{પ્રવાસી} \gu{આવ્યા}.} — \emph{gāmmā̃ ghaṇā pravāsī āvyā} — \textquotedblleft{}Many tourists came to the village.\textquotedblright{}
\end{quote}

Built on \textbf{\gu{પ્રવાસ}}, \emph{pravās}, 'trip'.

\begin{cousinweb}
\textbf{\gu{પ્રવાસી}} (\emph{pravāsī}) is Sanskrit \emph{pravāsin}, \textquotedblleft{}one who is away from home\textquotedblright{}, built on \emph{pravāsa}, \textquotedblleft{}a journey\textquotedblright{}, and taken over as a learned loan, a tatsama. CDIAL has no entry for it.
\end{cousinweb}

\begin{grammarlens}[title={What you've built}]
One more word for everyday life at A2.
\end{grammarlens}

\subsection*{Your turn}
\begin{itemize}
\raggedright
  \item \emph{Say it:} \emph{pravāsī}
  \item \emph{Say it:} \emph{pravāsī}, once more
  \item \emph{Say it:} \emph{gāmmā̃ ghaṇā pravāsī āvyā} — Many tourists came to the village.
  \item \emph{Recall:} say the Gujarati for to arrive, then the Gujarati for to get tired, then say \emph{pravāsī} again
\end{itemize}

\begin{tcolorbox}[breakable,colback=teal!4,colframe=teal!35!black,title={Before you move on}]
What does \gu{પ્રવાસી} mean? (\textquotedblleft{}Tourist, traveller\textquotedblright{}.) Say it once more.
\end{tcolorbox}
\section[\texorpdfstring{\gu{મુકામ} (mukām)}{મુકામ (mukām)}]{\gu{મુકામ} (mukām) — stopping place, destination}

\label{lesson:GU-C206-mukam}



\begin{quote}
Before the new one: say the Gujarati for to get tired, then the Gujarati for tourist.
\end{quote}

\subsection*{You'll want to know: \gu{મુકામ}}
\textbf{\gu{મુકામ}} — \emph{mukām} — \textquotedblleft{}stopping place, destination\textquotedblright{}.

\begin{quote}
\textbf{\gu{અમે} \gu{સાંજે} \gu{મુકામ} \gu{પર} \gu{પહોંચ્યા}.} — \emph{ame sā̃je mukām par pahõcyā} — \textquotedblleft{}We reached our stopping place in the evening.\textquotedblright{}
\end{quote}

\begin{cousinweb}
\textbf{\gu{મુકામ}} (\emph{mukām}) is a loan from Arabic \emph{muqām}, \emph{maqām}, \textquotedblleft{}a standing place, a place of stay\textquotedblright{}, which reached Gujarati through Persian. Its Arabic root is \emph{q-w-m}, \textquotedblleft{}to stand\textquotedblright{} (Lane).
\end{cousinweb}

\begin{grammarlens}[title={What you've built}]
One more word for everyday life at A2.
\end{grammarlens}

\subsection*{Your turn}
\begin{itemize}
\raggedright
  \item \emph{Say it:} \emph{mukām}
  \item \emph{Say it:} \emph{mukām}, once more
  \item \emph{Say it:} \emph{ame sā̃je mukām par pahõcyā} — We reached our stopping place in the evening.
  \item \emph{Recall:} say the Gujarati for to get tired, then the Gujarati for tourist, then say \emph{mukām} again
\end{itemize}

\begin{tcolorbox}[breakable,colback=teal!4,colframe=teal!35!black,title={Before you move on}]
What does \gu{મુકામ} mean? (\textquotedblleft{}Stopping place, destination\textquotedblright{}.) Say it once more.
\end{tcolorbox}
\section[\texorpdfstring{\gu{ગતિ} (gati)}{ગતિ (gati)}]{\gu{ગતિ} (gati) — speed}

\label{lesson:GU-C206-gati}



\begin{quote}
Before the new one: say the Gujarati for tourist, then the Gujarati for stopping place.
\end{quote}

\subsection*{You'll want to know: \gu{ગતિ}}
\textbf{\gu{ગતિ}} — \emph{gati} — \textquotedblleft{}speed\textquotedblright{}.

\begin{quote}
\textbf{\gu{ટ્રેનની} \gu{ગતિ} \gu{ધીમી} \gu{હતી}.} — \emph{ṭrennī gati dhīmī hatī} — \textquotedblleft{}The train was going slowly.\textquotedblright{}
\end{quote}

\begin{cousinweb}
\textbf{\gu{ગતિ}} (\emph{gati}) is Sanskrit \emph{gáti}, \textquotedblleft{}going, gait\textquotedblright{}, taken over whole as a learned loan, a tatsama (CDIAL 4009).
\end{cousinweb}

\begin{grammarlens}[title={What you've built}]
One more word for everyday life at A2.
\end{grammarlens}

\subsection*{Your turn}
\begin{itemize}
\raggedright
  \item \emph{Say it:} \emph{gati}
  \item \emph{Say it:} \emph{gati}, once more
  \item \emph{Say it:} \emph{ṭrennī gati dhīmī hatī} — The train was going slowly.
  \item \emph{Recall:} say the Gujarati for tourist, then the Gujarati for stopping place, then say \emph{gati} again
\end{itemize}

\begin{tcolorbox}[breakable,colback=teal!4,colframe=teal!35!black,title={Before you move on}]
What does \gu{ગતિ} mean? (\textquotedblleft{}Speed\textquotedblright{}.) Say it once more.
\end{tcolorbox}
\section[(dialogue)]{Words from chapters 201–206 — a review}

\label{lesson:GU-R206-review-201-206}



\begin{quote}
Say the Gujarati for to mend and cobbler.
\end{quote}

\subsection*{Your turn}
Say these aloud:

\begin{itemize}
\raggedright
  \item to mend, cobbler, washerman, barber, carpenter
  \item to steal, police station, to rob, to be found, to take out
  \item trip, luggage, to set out, to drop off, to see off
  \item timetable, enquiry, platform, to start off, to leave
  \item railway, seat, passenger, to halt, to get off
  \item to arrive, to get tired, tourist, stopping place, speed
\end{itemize}

\begin{tcolorbox}[breakable,colback=teal!4,colframe=teal!35!black,title={Before you move on}]
To mend? (\textbf{\gu{સુધારવું}}.) Cobbler? (\textbf{\gu{મોચી}}.) Washerman? (\textbf{\gu{ધોબી}}.) Barber? (\textbf{\gu{હજામ}}.) Carpenter? (\textbf{\gu{સુથાર}}.) To steal? (\textbf{\gu{ચોરવું}}.) Police station? (\textbf{\gu{થાણું}}.) To rob? (\textbf{\gu{લૂંટવું}}.) To be found? (\textbf{\gu{જડવું}}.) To take out? (\textbf{\gu{કાઢવું}}.) Trip? (\textbf{\gu{પ્રવાસ}}.) Luggage? (\textbf{\gu{સામાન}}.) To set out? (\textbf{\gu{નીકળવું}}.) To drop off? (\textbf{\gu{છોડવું}}.) To see off? (\textbf{\gu{વળાવવું}}.) Timetable? (\textbf{\gu{સમયપત્રક}}.) Enquiry? (\textbf{\gu{પૂછપરછ}}.) Platform? (\textbf{\gu{પ્લેટફોર્મ}}.) To start off? (\textbf{\gu{ઊપડવું}}.) To leave? (\textbf{\gu{છૂટવું}}.) Railway? (\textbf{\gu{રેલવે}}.) Seat? (\textbf{\gu{જગ્યા}}.) Passenger? (\textbf{\gu{ઉતારુ}}.) To halt? (\textbf{\gu{થોભવું}}.) To get off? (\textbf{\gu{ઊતરવું}}.) To arrive? (\textbf{\gu{પહોંચવું}}.) To get tired? (\textbf{\gu{થાકવું}}.) Tourist? (\textbf{\gu{પ્રવાસી}}.) Stopping place? (\textbf{\gu{મુકામ}}.) Speed? (\textbf{\gu{ગતિ}}.)
\end{tcolorbox}
